\chapter*{Abstract}
\addcontentsline{toc}{chapter}{Abstract}

Before a quantum gate is used in a computation or a sensing protocol it is characterized, and converse bounds on discriminating or estimating it, which limit every strategy including adaptive ones, are then evaluated on the characterized model. Is their sensitivity to characterization error a defect of the bounds or a property of the tasks they bound? The book opens with this question. Its first chapter shows that the supports of the Choi matrices separate genuine fragility, where the operational exponent itself jumps, from spurious fragility, where only single-letter converses jump. Near a noisy gate the sensitivity of the converses is governed by the noise floor, the smallest eigenvalue of the normalized Choi matrix: at first order it belongs to the task, whereas at second order the converses can overshoot the adaptive Stein and Chernoff exponents by factors that grow as the noise decreases. Certifying a converse from tomographic data costs a number of channel uses set by the same floor.

The rest of the book develops the theory behind these results. Converse bounds limit how well quantum channels can be discriminated, estimated or used to transmit information, and they are evaluated on a model of the channels that is known only approximately. We study when such bounds, and the operational exponents they bound, depend continuously on the channels, and what sets the size of the error when they do.

The divergences built from Kubo--Ando means, among them the geometric fidelity, the geometric R\'enyi divergences and the Belavkin--Staszewski relative entropy, are functions of a single positive contraction: the Arveson--Radon--Nikodym derivative of one channel with respect to the sum of the pair. In this coordinate they are continuous with dimension-free moduli. In the diamond norm, however, they are continuous at a pair of completely positive maps exactly when the Choi matrix of the sum is invertible or the fidelity vanishes, so ideal unitary gates are points of discontinuity. Near pairs whose Choi matrices are bounded below by $\lambda$, the fidelity is Lipschitz with a relative constant of the sharp order $\lambda^{-1}$.

Divergences built from hypothesis tests behave differently. The Umegaki relative entropy and the hockey-stick $f$-divergences are functionals of the testing region of a pair, and those with finite boundary values are uniformly continuous in every dimension. For channels, the adaptive Stein exponent is continuous exactly when the Choi matrix of the second channel is invertible or the exponent is infinite. We also study when the hockey-stick R\'enyi quasi-entropies are super-multiplicative under tensor products, and prove several cases of the corresponding conjecture.

For adaptive discrimination, every protocol is a controlled random walk on the spectrum of the derivative, which gives a geometric Chernoff converse with an explicit second-order term. Single-letter converses of this kind can be infinite where the adaptive exponent tends to zero, as for amplitude-damping channels with nearby damping rates. A chapter of proofs then returns to noisy gates and establishes the results of the first chapter, and a final chapter collects the open problems.
